%智能体系统近期广泛应用->MAS->利用tool-like-agent 是新的解法->
%问题：人工创建/静态上下文->

\begin{abstract}
% As open-world, long-horizon tasks become increasingly complex, Multi-Agent Systems (MAS) have been a common solution. More recently, a \emph{tool-like agent} paradigm has emerged, where sub-agents are invoked as modular executors to improve composability and reduce coordination burden. However, today’s designs remain rigid: sub-agents are typically implemented either as context-isolated threads without specialized capabilities or as static roles that cannot adapt to different tasks. This stems from the lack of a \emph{dynamic abstraction} view, forcing systems into these two limited design patterns.

% As tasks become increasingly complex and long-horizon, a subagent-as-tools paradigm has emerged to xx. However, we find existing designs do not treat these agents from a \emph{dynamic} view, leading to poor adaptibility: sub-agents are typically implemented either as context-isolated threads that are not specialized or as static roles that require human-engineering.

Language agents have shown strong promise for task automation. 
% 
Realizing this promise for increasingly complex, long-horizon tasks has driven the rise of a sub-agent-as-tools paradigm for multi-turn task solving. 
% 
However, existing designs still lack a \emph{dynamic abstraction} view of sub-agents, thereby hurting adaptability.
% 
We address this challenge with a unified, framework-agnostic agent abstraction that models any agent as a tuple $\langle Instruction, Context,Tools,Model \rangle$. This tuple acts as a compositional recipe for capabilities, enabling the system to spawn specialized executors for each task on demand. 
Building on this abstraction, we introduce an agentic system \textsc{AOrchestra}, where the central orchestrator concretizes the tuple at each step: it curates task-relevant context, selects tools and models, and delegates execution via on-the-fly automatic agent creation.
Such designs enable reducing human engineering efforts, and remain framework-agnostic with plug-and-play support for diverse agents as task executors. It also enables a controllable performance–cost trade-off, allowing the system to approach Pareto-efficient.
Across three challenging benchmarks (GAIA, SWE-Bench, Terminal-Bench), \textsc{AOrchestra} achieves 16.28\% relative improvement against the strongest baseline when paired with Gemini-3-Flash. The code is available at: \url{https://github.com/FoundationAgents/AOrchestra  }

% As open-world, long-horizon interactive tasks grow increasingly complex, the scalability and adaptability of agentic systems in dynamic environments remain challenging. 
% Existing agentic systems often rely on rigid context-sharing strategies and static, human-engineered sub-agents. 
% In this paper, we propose a \textbf{lightweight orchestration framework} grounded in a unified abstraction that represents any agent as a four-tuple \textbf{$\langle \textsc{LM}, \textsc{Prompt}, \textsc{Tools}, \textsc{Context} \rangle$}. Building on this abstraction, we introduce a central \textbf{\textsc{Orchestrator}} that decomposes tasks, schedules and terminates subtasks, and composes \textbf{plug-and-play sub-agents} on demand for each subtask. 
% This design increases end-to-end autonomy by reducing manual engineering and enabling modular, lightweight composition, and it delivers \textbf{consistent gains across diverse agent frameworks} through its plug-in orchestration layer. 
% We demonstrate that our framework improves task execution, scalability, and adaptability, addressing key limitations of existing agentic systems while supporting more autonomous operation in complex environments.
\end{abstract}
